\documentclass{article}
\usepackage[T1]{fontenc}
\usepackage{lmodern}
\usepackage{graphicx}
\usepackage{amsmath,amssymb}
\usepackage[a4paper,margin=2cm]{geometry}
\usepackage[absolute,overlay]{textpos}
\setlength{\TPHorizModule}{1mm}
\setlength{\TPVertModule}{1mm}
\usepackage[indonesian]{babel}
\usepackage{hyperref}
\setlength{\emergencystretch}{2em}
% unit: Halaman 8

\begin{document}

% === Halaman 8 (python/native) ===

• Jika plainteks/cipherteks hanya menggunakan 26 huruf alfabet (A–Z), maka entropi 

Shannon maksimum terjadi ketika seluruh 26 huruf memiliki probabilitas yang sama:

                      𝑝𝑖= 

1

26

   sehingga 


𝐻𝑚𝑎𝑥=  −σ𝑖=1

26 𝑝𝑖 log2 𝑝𝑖 


= −σ𝑖=1

26

1

26 log2

1

26 = log2 26 = 4.7004 bit/karakter

   Jadi, untuk teks yang hanya terdiri dari 26 huruf alfabet, entropi Shannon maksimum  

adalah sekitar 4,7004 bit/karakter.

• Jika pesan terdiri dari 256 karakter ASCII, maka entropi maksimum adalah log2 256 = 8 

• Makin besar entropi, makin sulit memecahkan cipherteks. 

8

\end{document}
